\chapter{把很多数整理成三个块，同时留下余数差异}\label{sec:compression}
\section{位置不够了，换一个工作目标}
如果 $n=2p+1$，两个 $p$ 块已经用掉 $2p$ 个位置，只剩一个单点。前面那种“把两块都留住，再到外面找一组”的方法没法继续。

现在我们允许平均产生分数。目标是把所有位置整理成
\[(A^p,B^p,C^r),\qquad n=2p+r.
\]
这还没有解决问题，但复杂度下降了：三种数值受一个零和方程约束，只剩两个自由参数。我们可以集中研究这两个参数怎样变化。

难点是不能随意分组。三平均的 $(3^5,1,-16)$ 已经告诉我们，选错一组会毁掉唯一的模七见证。因此，下面先把余数差异放在容易保护的位置上，再合并其他位置。

\section{同一个整数，要同时躲开几个禁区}
两块互补混合：第一组选 $p-k$ 份 $a$、$k$ 份 $b$，另一组选剩余的互补份数。结果是
\begin{equation}\label{eq:mix-pair}
a_k=a+\frac{k}{p}(b-a),\qquad b_k=b+\frac{k}{p}(a-b).
\end{equation}
固定参照值 $w$。模某个危险素数 $\ell$，$a_k-w$ 是 $k$ 的一次式。只要 $a-w,b-w$ 不同时为零，这个一次式要么一直非零，要么只禁止一个 $k$ 余数。

这时中国剩余定理只能告诉我们某个整数存在，却不保证它在 $1,\ldots,p-1$ 里。实际的 $k$ 是取用份数，不能无限大。因此还要用 $n<4p$ 提供的素因子数量限制，证明这个短区间中就有一个同时可用的 $k$。
\InsertProof{lem:short-choice}
证明中的 $\rho=\varphi(M)/M$ 可以先理解成“模 $M$ 的整数里，与 $M$ 互素的比例”。容斥给出主项 $L\rho$，再减去一个可控制的边界误差。如果主项比误差大，至少剩一个可用整数。

\section{把分散的差异集中到一对数上}
例如某块 $a$ 与单点 $w$ 已在模七不同，但模十一还相同；另一块 $b$ 在模十一与 $w$ 不同。用互补混合，选一个同时避开两个坏配比的 $k$，就能使同一个 $a_k-w$ 在两个素数处都非零。

如果差异在一个单点 $x$ 上，先平均 $(b^{p-1},x)$。留下的一份旧 $b$ 接替单点角色，新值接替 $b$ 块。对应二元变换为
\[(b,x)\mapsto\left(\frac{(p-1)b+x}{p},b\right).
\]
它模危险素数可逆，又固定常值配置，所以不会把原有的整体差异全部抹掉。
\InsertProof{lem:concentrate}

\section{最紧的情况，怎样先造出两个块}
还不能默认我们一开始已有两个块。先保护各素数的不同余位置，造出一个块。剩下是多个单点。对任一单点，刚才的块--单点交换都可以执行。

接下来的交换子公式看起来突然，其目的很简单：让一个目标单点的差值增加或减少另一个单点差值的倍数，并让其余多数差值保持不动。它就是高维坐标中的一次“加倍数”操作。

此时只需要在模 $\rad(n)$ 下得到这个作用。有限可逆矩阵的逆等于它的某个非负次幂，所以公式里的逆可以通过多做正向交换得到相同的余数作用。这不把实际有理数值倒回去；入口阶段只需要余数正确。
\InsertProof{lem:two-block-entry}
最关键的容量配合是：交换子失效的素数都整除 $r+2$，其个数不超过剩余单点数 $r$。失效处用单点保存，其他处用交换子搬运。因此即使 $n$ 有很多素因子，也不需要为每个素因子永久占据一个额外位置。

\section{为什么核心要写成这个样子}
设三块零和：$pA+pB+rC=0$，且 $\gcd(p,r)=1$。在整数代表中 $p\mid C$，写 $C=-pz$，令 $A=a$，便得到
\[B=rz-a.
\]
所以
\[K_{p,r}(a,z)=(a^p,(rz-a)^p,(-pz)^r)
\]
只是零和方程的一个方便参数化。令 $y=a+pz=A-C$，另外的块差可写成 $nz-2y$。本原性给 $\gcd(y,z)=1$，于是
\[G=\gcd(y,n).
\]
我们将来要控制的“第二坐标是模 $n$ 单位”，正是最初的 $G=1$ 条件。
\InsertProof{prop:entry}

\section{上半带多出来的一个块，怎样换成大载体}
当 $n=3p+s$ 时，最容易先得到的是 $(2p,p,s)$ 重数：三个完整 $p$ 块先折叠出一个 $2p$ 块。后面的共同矩阵工具则更喜欢两个等大的 $p$ 块，因此要把它换成 $(p,p,p+s)$。

这里不用任何抽象“拆块定理”。先平均一组，从尚未用掉的旧 $A$ 中留下 $p+s$ 份，再平均其余 $p$ 份，一共两次。唯一需要选好的，是第一组从三种值各取多少。
\InsertProof{lem:shape-change}
到这里，任意合法输入已经被送入一个覆盖全部位置的二维参数族。后面用矩阵工作时，别把两个参数误看成只剩两个位置；原来的 $n$ 个位置一个也没少。
